\makeatletter
\def\input@path{{./}{../}{../../}{../../../}{../../../../}{preamble/}{../preamble/}{../../preamble/}{../../../preamble/}{../../../../preamble/}{book2026-2027/preamble/}{book2026-2027/}}
\makeatother

\input{preamble_G10_T1_C1C2C3}
\input{preamble_exam_layout}
\input{graphs/uniform_linear_probability_distributions.tex}
\input{graphs/probability_distribution_graphs.tex}
\input{graphs/normal_distribution_graphs.tex}

\title{Grade 10 -- Term 1 \\ C1--C2 Classroom Activity with Solutions -- Version 2}
\author{}
\date{}

\pagestyle{empty}

\geometry{
	margin=0.3in,
	bottom=0.3in
}

% Flexible table columns used on the solution pages.  The Y column from the
% shared exam preamble has a fixed width, so four of them cannot fit on a line.

\title{\vspace{-15mm} \textbf{Global Economics}\\
	Grade 10 -- Term 1 -- C1-C2 Classroom Activity -- Version 2}
\author{}
\date{}
\pagestyle{empty}

\hypersetup{pdftitle={Grade 10 -- Term 1 -- C1-C2 Classroom Activity -- Version 2}}
\begin{document}

	\maketitle\thispagestyle{empty}
	\vspace{-22mm}

\textbf{(C1)} \textit{Characterises decision-making problems under uncertainty by identifying decision alternatives, states of nature, consequences, and the economic or financial context of the decision.}. 

\textbf{(C2)} \textit{Organises decision problems using payoff tables, placing alternatives, states of nature, and payoffs in a clear structure for comparison.}. \vspace{-3mm}
\section*{Context Descriptions}
\vspace{-3mm}
\textit{All per-unit contributions and adjustments are already net of relevant per-unit costs. Fixed amounts apply once over the stated activity period. All state adjustments apply to every alternative. No other monetary components are required.}
\begin{enumerate}[leftmargin=0.22in]
\item
A textbook publisher must choose a digital licence or a print licence before the education board announces whether it will retain or revise its syllabus. These are external syllabus decisions. Analysts have already calculated the complete profit for each outcome, in thousand USD. The digital licence yields 23 thousand USD if the syllabus is retained and 14 thousand USD if it is revised. The print licence yields 34 thousand USD if the syllabus is retained and 9 thousand USD if it is revised.

\item
A rehearsal studio must choose a daytime contract or an evening contract before an independent organiser announces an extended or shortened festival. The daytime contract provides a fixed net contribution of 640 thousand JPY; the evening contract provides a fixed net contribution of 480 thousand JPY. An extended festival adds a fixed access payment of 210 thousand JPY to either contract. A shortened festival instead imposes a fixed cancellation deduction of 120 thousand JPY on either contract. No output quantity enters these fixed payments.

\item
An archive contractor must select a flatbed scanner or an overhead camera before ambient humidity during a project is known. The external conditions will be dry air or humid air. The flatbed scanner contributes 7 thousand CNY per completed archive collection; the overhead camera contributes 9 thousand CNY per completed collection. With the flatbed scanner, 13 collections would be completed in dry air and 9 in humid air. With the overhead camera, 11 would be completed in dry air and 7 in humid air. Dry air triggers a fixed preservation-contract bonus of 8 thousand CNY; humid air causes a fixed climate-control loss of 7 thousand CNY. Each fixed adjustment applies to either equipment choice.

\item
A ceramics manufacturer must choose a tableware line or a tile line before external wholesale orders are known. Orders will be brisk or quiet. The tableware line contributes 5 thousand EUR per completed firing batch and a fixed partner contribution of 9 thousand EUR. The tile line contributes 7 thousand EUR per completed firing batch and a fixed partner contribution of 13 thousand EUR. Brisk orders add 4 thousand EUR per batch to either line; quiet orders add 2 thousand EUR per batch to either line. The tableware line would complete 18 batches with brisk orders and 11 with quiet orders. The tile line would complete 15 and 8 batches under those respective conditions.

\item
A vocational training centre must choose a woodworking programme or a metalworking programme before the regional authority announces increased, baseline, or reduced training funding. The woodworking programme contributes 6 thousand GBP per completed course and a fixed partner payment of 9 thousand GBP. The metalworking programme contributes 7 thousand GBP per completed course and a fixed partner payment of 11 thousand GBP. For either programme, increased funding adds 3 thousand GBP per course and a fixed administration grant of 8 thousand GBP. Baseline funding adds 2 thousand GBP per course and has no fixed state adjustment. Reduced funding deducts 2 thousand GBP per course and imposes a fixed administration loss of 6 thousand GBP. Under increased, baseline, and reduced funding respectively, the centre would complete 12, 8, and 3 woodworking courses, or 16, 11, and 6 metalworking courses. These are the complete conditional course quantities.

\end{enumerate}
\newpage
\newpage
\section{Solutions: Extraction -- Version 2}
\phantomsection
\label{sec:solutions-extraction}
\label{toc}
\SectionProblemsTableOfContents{
\SectionProblemLink{sec:problem1-textbook-publishing-rights-c1-c2}{Textbook publishing rights (C1--C2)} \\
\SectionProblemLink{sec:problem2-rehearsal-studio-contract-c1-c2}{Rehearsal studio contract (C1--C2)} \\
\SectionProblemLink{sec:problem3-archive-imaging-equipment-c1-c2}{Archive imaging equipment (C1--C2)} \\
\SectionProblemLink{sec:problem4-ceramics-production-line-c1-c2}{Ceramics production line (C1--C2)} \\
\SectionProblemLink{sec:problem5-vocational-training-programme-c1-c2}{Vocational training programme (C1--C2)} \\
}
\subsection{Exercise 1: Textbook publishing rights}
\phantomsection
\label{sec:problem1-textbook-publishing-rights-c1-c2}
\paragraph*{C1 - Interpreting decision alternatives, events, consequences, and states.}
\phantomsection
\label{subsec:problem1-textbook-publishing-rights-c1}
\DecisionElementsTable
{Digital licence; Print licence}
{Syllabus retained; Syllabus revised}
{Whether the education board retains or revises its syllabus after the publishing licence is chosen}
{Profit in thousand USD for each alternative--state pair. All four profits are supplied directly.}
\secInicio[sec:solutions-extraction][sec:problem1-textbook-publishing-rights-c1-c2]
\paragraph*{C2 - Building the payoff table.}
\leavevmode\par
\phantomsection
\label{subsec:problem1-textbook-publishing-rights-c2}
\textbf{Relevant values.} The context supplies the complete profit for every licence--syllabus pair.
\begin{center}
\small
\begin{tabularx}{\textwidth}{LLLL}
\toprule
Abstract component & Context component & Math symbol & Corresponding value \\
\midrule
Direct payoff & Complete profit for the chosen licence under the board's syllabus decision & $P_{ay}(A_i,S_j)$ & Digital--retained: 23 thousand USD; Digital--revised: 14 thousand USD; Print--retained: 34 thousand USD; Print--revised: 9 thousand USD \\
\bottomrule
\end{tabularx}
\end{center}
\textbf{Payoff equation.} $P_{ay}(A_i,S_j)=\text{direct profit stated for the licence--syllabus pair}$.
\textbf{Calculations.} No components need to be added: each symbolic payoff is assigned its stated numerical value.
\begin{center}
\begin{tabularx}{\textwidth}{L>{\centering\arraybackslash}X>{\centering\arraybackslash}X}
\toprule
Licence & Syllabus retained ($S_1$) & Syllabus revised ($S_2$) \\
\midrule
Digital licence ($A_1$) & $P_{ay}(A_1,S_1)=23$ & $P_{ay}(A_1,S_2)=14$ \\
Print licence ($A_2$) & $P_{ay}(A_2,S_1)=34$ & $P_{ay}(A_2,S_2)=9$ \\
\bottomrule
\end{tabularx}
\end{center}
\begingroup\renewcommand{\CalculationStateCount}{2}
\FinalPayoffTable{profit in thousand USD}{Syllabus retained ($S_1$) & Syllabus revised ($S_2$)}{
Digital licence ($A_1$) & 23 & 14 \\
Print licence ($A_2$) & 34 & 9 \\
}
\endgroup
\subsection{Exercise 2: Rehearsal studio contract}
\phantomsection
\label{sec:problem2-rehearsal-studio-contract-c1-c2}
\paragraph*{C1 - Interpreting decision alternatives, events, consequences, and states.}
\phantomsection
\label{subsec:problem2-rehearsal-studio-contract-c1}
\DecisionElementsTable
{Daytime contract; Evening contract}
{Extended festival; Shortened festival}
{The festival duration announced by the independent organiser after the studio contract is chosen}
{Profit in thousand JPY for each alternative--state pair. Components: alternative-specific fixed amount, state-specific fixed adjustment.}
\secInicio[sec:solutions-extraction][sec:problem2-rehearsal-studio-contract-c1-c2]
\begin{CalculationLandscape}{11in}{8.5in}
\paragraph*{C2 - Building the payoff table.}
\leavevmode\par
\phantomsection
\label{subsec:problem2-rehearsal-studio-contract-c2}
\SymbolsTable{
$A_i$ & $A_1$: Daytime contract; $A_2$: Evening contract & Alternative labels, in row order. \\
$S_j$ & $S_1$: Extended festival; $S_2$: Shortened festival & State labels, in column order; revealed after the choice. \\
$f^A(A_i)$ & $A_1$: 640 thousand JPY; $A_2$: 480 thousand JPY & Fixed net contribution from the chosen studio contract. \\
$f^S(S_j)$ & $S_1$: 210 thousand JPY; $S_2$: $-120$ thousand JPY & Fixed access payment or cancellation deduction, applied once to either contract. \\
}
\textbf{Payoff equation.}
	\[
		P_{ay}(A_i,S_j)=f^A(A_i)+f^S(S_j).
	\]
	\CalculationPair{Additive payoff}{
\begin{tabular}{@{}lcc@{}}
\toprule Alternative & $S_1$ & $S_2$ \\ \midrule
$A_1$ & \SymbolicCell{f^A(A_1) \\ {}+f^S(S_1)}{P_{ay}(A_1,S_1)} & \SymbolicCell{f^A(A_1) \\ {}+f^S(S_2)}{P_{ay}(A_1,S_2)} \\[5pt]
$A_2$ & \SymbolicCell{f^A(A_2) \\ {}+f^S(S_1)}{P_{ay}(A_2,S_1)} & \SymbolicCell{f^A(A_2) \\ {}+f^S(S_2)}{P_{ay}(A_2,S_2)} \\[5pt]
\bottomrule
\end{tabular}
}{
\begin{tabular}{@{}lcc@{}}
\toprule Alternative & $S_1$ & $S_2$ \\ \midrule
$A_1$ & \CalcResult{640 \\ {}+210}{850} & \CalcResult{640 \\ {}-120}{520} \\[5pt]
$A_2$ & \CalcResult{480 \\ {}+210}{690} & \CalcResult{480 \\ {}-120}{360} \\[5pt]
\bottomrule
\end{tabular}
}
Negative fixed amounts are deducted once. A negative per-unit adjustment reduces the contribution per completed unit; a negative payoff is a shortfall.
\begingroup\renewcommand{\CalculationStateCount}{2}
\FinalPayoffTable{profit in thousand JPY}{Extended festival ($S_1$) & Shortened festival ($S_2$)}{
Daytime contract ($A_1$) & 850 & 520 \\
Evening contract ($A_2$) & 690 & 360 \\
}
\endgroup
\end{CalculationLandscape}
\newpage
\subsection{Exercise 3: Archive imaging equipment}
\phantomsection
\label{sec:problem3-archive-imaging-equipment-c1-c2}
\paragraph*{C1 - Interpreting decision alternatives, events, consequences, and states.}
\phantomsection
\label{subsec:problem3-archive-imaging-equipment-c1}
\DecisionElementsTable
{Flatbed scanner; Overhead camera}
{Dry air; Humid air}
{Ambient humidity during the archive project, revealed after the imaging equipment is selected}
{Profit in thousand CNY for each alternative--state pair. Components: completed output quantity, alternative-specific net contribution per unit, state-specific fixed adjustment.}
\secInicio[sec:solutions-extraction][sec:problem3-archive-imaging-equipment-c1-c2]
\begin{CalculationLandscape}{11in}{8.5in}
\paragraph*{C2 - Building the payoff table.}
\leavevmode\par
\phantomsection
\label{subsec:problem3-archive-imaging-equipment-c2}
\SymbolsTable{
$A_i$ & $A_1$: Flatbed scanner; $A_2$: Overhead camera & Alternative labels, in row order. \\
$S_j$ & $S_1$: Dry air; $S_2$: Humid air & State labels, in column order; revealed after the choice. \\
$q(A_i,S_j)$ & $A_1$: 13, 9 collections; $A_2$: 11, 7 collections (state order) & Completed archive collections for each equipment--humidity pair. \\
$v^A(A_i)$ & $A_1$: 7 thousand CNY/collection; $A_2$: 9 thousand CNY/collection & Net contribution per completed collection from the chosen equipment. \\
$f^S(S_j)$ & $S_1$: 8 thousand CNY; $S_2$: $-7$ thousand CNY & Fixed preservation bonus or climate-control loss for either choice. \\
}
\textbf{Payoff equation.}
	\[
		P_{ay}(A_i,S_j)=q(A_i,S_j)v^A(A_i)+f^S(S_j).
	\]
	\CalculationPair{Additive payoff}{
\begin{tabular}{@{}lcc@{}}
\toprule Alternative & $S_1$ & $S_2$ \\ \midrule
$A_1$ & \SymbolicCell{q(A_1,S_1)\times v^A(A_1) \\ {}+f^S(S_1)}{P_{ay}(A_1,S_1)} & \SymbolicCell{q(A_1,S_2)\times v^A(A_1) \\ {}+f^S(S_2)}{P_{ay}(A_1,S_2)} \\[5pt]
$A_2$ & \SymbolicCell{q(A_2,S_1)\times v^A(A_2) \\ {}+f^S(S_1)}{P_{ay}(A_2,S_1)} & \SymbolicCell{q(A_2,S_2)\times v^A(A_2) \\ {}+f^S(S_2)}{P_{ay}(A_2,S_2)} \\[5pt]
\bottomrule
\end{tabular}
}{
\begin{tabular}{@{}lcc@{}}
\toprule Alternative & $S_1$ & $S_2$ \\ \midrule
$A_1$ & \CalcResult{13\times 7 \\ {}+8}{99} & \CalcResult{9\times 7 \\ {}-7}{56} \\[5pt]
$A_2$ & \CalcResult{11\times 9 \\ {}+8}{107} & \CalcResult{7\times 9 \\ {}-7}{56} \\[5pt]
\bottomrule
\end{tabular}
}
Negative fixed amounts are deducted once. A negative per-unit adjustment reduces the contribution per completed unit; a negative payoff is a shortfall.
\begingroup\renewcommand{\CalculationStateCount}{2}
\FinalPayoffTable{profit in thousand CNY}{Dry air ($S_1$) & Humid air ($S_2$)}{
Flatbed scanner ($A_1$) & 99 & 56 \\
Overhead camera ($A_2$) & 107 & 56 \\
}
\endgroup
\end{CalculationLandscape}
\newpage
\subsection{Exercise 4: Ceramics production line}
\phantomsection
\label{sec:problem4-ceramics-production-line-c1-c2}
\paragraph*{C1 - Interpreting decision alternatives, events, consequences, and states.}
\phantomsection
\label{subsec:problem4-ceramics-production-line-c1}
\DecisionElementsTable
{Tableware line; Tile line}
{Brisk wholesale orders; Quiet wholesale orders}
{The level of wholesale orders received after the ceramics production line is selected}
{Profit in thousand EUR for each alternative--state pair. Components: completed output quantity, alternative-specific net contribution per unit, state-specific net adjustment per unit, alternative-specific fixed amount.}
\secInicio[sec:solutions-extraction][sec:problem4-ceramics-production-line-c1-c2]
\begin{CalculationLandscape}{11in}{8.5in}
\paragraph*{C2 - Building the payoff table.}
\leavevmode\par
\phantomsection
\label{subsec:problem4-ceramics-production-line-c2}
\SymbolsTable{
$A_i$ & $A_1$: Tableware line; $A_2$: Tile line & Alternative labels, in row order. \\
$S_j$ & $S_1$: Brisk wholesale orders; $S_2$: Quiet wholesale orders & State labels, in column order; revealed after the choice. \\
$q(A_i,S_j)$ & $A_1$: 18, 11 batches; $A_2$: 15, 8 batches (state order) & Completed firing batches for each line--order pair. \\
$v^A(A_i)$ & $A_1$: 5 thousand EUR/batch; $A_2$: 7 thousand EUR/batch & Net contribution per firing batch from the chosen line. \\
$v^S(S_j)$ & $S_1$: 4 thousand EUR/batch; $S_2$: 2 thousand EUR/batch & Additional net contribution per batch under either line. \\
$f^A(A_i)$ & $A_1$: 9 thousand EUR; $A_2$: 13 thousand EUR & Positive fixed partner contribution from the chosen production line. \\
}
\textbf{Payoff equation.}
	\[
		P_{ay}(A_i,S_j)
		=q(A_i,S_j)\left[v^A(A_i)+v^S(S_j)\right]+f^A(A_i).
	\]
	\CalculationPair{Additive payoff}{
\begin{tabular}{@{}lcc@{}}
\toprule Alternative & $S_1$ & $S_2$ \\ \midrule
$A_1$ & \SymbolicCell{q(A_1,S_1)\times \\ [v^A(A_1)+v^S(S_1)] \\ {}+f^A(A_1)}{P_{ay}(A_1,S_1)} & \SymbolicCell{q(A_1,S_2)\times \\ [v^A(A_1)+v^S(S_2)] \\ {}+f^A(A_1)}{P_{ay}(A_1,S_2)} \\[5pt]
$A_2$ & \SymbolicCell{q(A_2,S_1)\times \\ [v^A(A_2)+v^S(S_1)] \\ {}+f^A(A_2)}{P_{ay}(A_2,S_1)} & \SymbolicCell{q(A_2,S_2)\times \\ [v^A(A_2)+v^S(S_2)] \\ {}+f^A(A_2)}{P_{ay}(A_2,S_2)} \\[5pt]
\bottomrule
\end{tabular}
}{
\begin{tabular}{@{}lcc@{}}
\toprule Alternative & $S_1$ & $S_2$ \\ \midrule
$A_1$ & \CalcResult{18\times \\ [5+4] \\ {}+9}{171} & \CalcResult{11\times \\ [5+2] \\ {}+9}{86} \\[5pt]
$A_2$ & \CalcResult{15\times \\ [7+4] \\ {}+13}{178} & \CalcResult{8\times \\ [7+2] \\ {}+13}{85} \\[5pt]
\bottomrule
\end{tabular}
}
\begingroup\renewcommand{\CalculationStateCount}{2}
\FinalPayoffTable{profit in thousand EUR}{Brisk wholesale orders ($S_1$) & Quiet wholesale orders ($S_2$)}{
Tableware line ($A_1$) & 171 & 86 \\
Tile line ($A_2$) & 178 & 85 \\
}
\endgroup
\end{CalculationLandscape}
\subsection{Exercise 5: Vocational training programme}
\phantomsection
\label{sec:problem5-vocational-training-programme-c1-c2}
\paragraph*{C1 - Interpreting decision alternatives, events, consequences, and states.}
\phantomsection
\label{subsec:problem5-vocational-training-programme-c1}
\DecisionElementsTable
{Woodworking programme; Metalworking programme}
{Increased funding; Baseline funding; Reduced funding}
{The regional authority's training-funding level, announced after the programme is selected}
{Profit in thousand GBP for each alternative--state pair. Components: completed output quantity, alternative-specific net contribution per unit, state-specific net adjustment per unit, alternative-specific fixed amount, state-specific fixed adjustment.}
\secInicio[sec:solutions-extraction][sec:problem5-vocational-training-programme-c1-c2]
\begin{CalculationLandscape}{11in}{8.5in}
\paragraph*{C2 - Building the payoff table.}
\leavevmode\par
\phantomsection
\label{subsec:problem5-vocational-training-programme-c2}
\SymbolsTable{
$A_i$ & $A_1$: Woodworking programme; $A_2$: Metalworking programme & Alternative labels, in row order. \\
$S_j$ & $S_1$: Increased funding; $S_2$: Baseline funding; $S_3$: Reduced funding & State labels, in column order; revealed after the choice. \\
$q(A_i,S_j)$ & $A_1$: 12, 8, 3 courses; $A_2$: 16, 11, 6 courses (state order) & Completed courses for each programme--funding pair. \\
$v^A(A_i)$ & $A_1$: 6 thousand GBP/course; $A_2$: 7 thousand GBP/course & Net contribution per completed course from the chosen programme. \\
$v^S(S_j)$ & $S_1$: 3 thousand GBP/course; $S_2$: 2 thousand GBP/course; $S_3$: $-2$ thousand GBP/course & Per-course funding addition or deduction, for either programme. \\
$f^A(A_i)$ & $A_1$: 9 thousand GBP; $A_2$: 11 thousand GBP & Positive fixed partner payment for the chosen programme. \\
$f^S(S_j)$ & $S_1$: 8 thousand GBP; $S_2$: 0 thousand GBP; $S_3$: $-6$ thousand GBP & Fixed administration grant, zero adjustment, or administration loss for either programme. \\
}
\textbf{Payoff equation.}
	\[
		P_{ay}(A_i,S_j)
		=q(A_i,S_j)\left[v^A(A_i)+v^S(S_j)\right]+f^A(A_i)+f^S(S_j).
	\]
	\newpage
\CalculationPair{Additive payoff}{
\begin{tabular}{@{}lccc@{}}
\toprule Alternative & $S_1$ & $S_2$ & $S_3$ \\ \midrule
$A_1$ & \SymbolicCell{q(A_1,S_1)\times \\ [v^A(A_1)+v^S(S_1)] \\ {}+f^A(A_1) \\ {}+f^S(S_1)}{P_{ay}(A_1,S_1)} & \SymbolicCell{q(A_1,S_2)\times \\ [v^A(A_1)+v^S(S_2)] \\ {}+f^A(A_1) \\ {}+f^S(S_2)}{P_{ay}(A_1,S_2)} & \SymbolicCell{q(A_1,S_3)\times \\ [v^A(A_1)+v^S(S_3)] \\ {}+f^A(A_1) \\ {}+f^S(S_3)}{P_{ay}(A_1,S_3)} \\[5pt]
$A_2$ & \SymbolicCell{q(A_2,S_1)\times \\ [v^A(A_2)+v^S(S_1)] \\ {}+f^A(A_2) \\ {}+f^S(S_1)}{P_{ay}(A_2,S_1)} & \SymbolicCell{q(A_2,S_2)\times \\ [v^A(A_2)+v^S(S_2)] \\ {}+f^A(A_2) \\ {}+f^S(S_2)}{P_{ay}(A_2,S_2)} & \SymbolicCell{q(A_2,S_3)\times \\ [v^A(A_2)+v^S(S_3)] \\ {}+f^A(A_2) \\ {}+f^S(S_3)}{P_{ay}(A_2,S_3)} \\[5pt]
\bottomrule
\end{tabular}
}{
\begin{tabular}{@{}lccc@{}}
\toprule Alternative & $S_1$ & $S_2$ & $S_3$ \\ \midrule
$A_1$ & \CalcResult{12\times \\ [6+3] \\ {}+9 \\ {}+8}{125} & \CalcResult{8\times \\ [6+2] \\ {}+9 \\ {}+0}{73} & \CalcResult{3\times \\ [6-2] \\ {}+9 \\ {}-6}{15} \\[5pt]
$A_2$ & \CalcResult{16\times \\ [7+3] \\ {}+11 \\ {}+8}{179} & \CalcResult{11\times \\ [7+2] \\ {}+11 \\ {}+0}{110} & \CalcResult{6\times \\ [7-2] \\ {}+11 \\ {}-6}{35} \\[5pt]
\bottomrule
\end{tabular}
}
Negative fixed amounts are deducted once. A negative per-unit adjustment reduces the contribution per completed unit; a negative payoff is a shortfall.
\begingroup\renewcommand{\CalculationStateCount}{3}
\FinalPayoffTable{profit in thousand GBP}{Increased funding ($S_1$) & Baseline funding ($S_2$) & Reduced funding ($S_3$)}{
Woodworking programme ($A_1$) & 125 & 73 & 15 \\
Metalworking programme ($A_2$) & 179 & 110 & 35 \\
}
\endgroup
\end{CalculationLandscape}
\end{document}
